\documentclass[12pt,a4paper]{article}

%% ---------- mathematics ----------
\usepackage{amsmath,amsthm,amssymb,mathtools}

%% ---------- graphics & diagrams ----------
\usepackage{tikz}
\usetikzlibrary{arrows.meta,positioning,shapes.geometric,calc,fit,backgrounds,matrix,decorations.pathreplacing}
\usepackage{graphicx}

%% ---------- tables ----------
\usepackage{booktabs,tabularx,multirow,longtable,array}

%% ---------- algorithms ----------
\usepackage[ruled,vlined,linesnumbered]{algorithm2e}

%% ---------- coloured boxes ----------
\usepackage[most]{tcolorbox}

%% ---------- hyperlinks & cross-refs ----------
\usepackage{hyperref}
\usepackage[capitalize,noabbrev]{cleveref}

%% ---------- page layout ----------
\usepackage[margin=1in]{geometry}
\usepackage{fancyhdr}
\usepackage{enumitem}
\usepackage{xcolor}

%% ---------- fonts & misc ----------
\usepackage{microtype}

%% ============================================================
%%  Custom tcolorbox environments
%% ============================================================
\newtcolorbox{keyinsight}[1][]{%
  colback=blue!5!white,colframe=blue!75!black,
  fonttitle=\bfseries,title=Key Insight,#1}

\newtcolorbox{example}[1][]{%
  colback=green!5!white,colframe=green!75!black,
  fonttitle=\bfseries,title=Example,#1}

\newtcolorbox{warningbox}[1][]{%
  colback=red!5!white,colframe=red!75!black,
  fonttitle=\bfseries,title=Warning,#1}

%% ============================================================
%%  amsthm environments
%% ============================================================
\newtheorem{theorem}{Theorem}[section]
\newtheorem{lemma}[theorem]{Lemma}
\newtheorem{definition}{Definition}[section]
\newtheorem{remark}{Remark}[section]
\newtheorem{proposition}[theorem]{Proposition}
\newtheorem{corollary}[theorem]{Corollary}

%% ============================================================
%%  Custom commands (collected from all sessions in this file)
%% ============================================================
\newcommand{\xt}{x_t}
\newcommand{\ht}{h_t}
\newcommand{\yt}{y_t}
\newcommand{\Whh}{W_{hh}}
\newcommand{\Wxh}{W_{xh}}
\newcommand{\Why}{W_{hy}}
\newcommand{\htmo}{h_{t-1}}
\newcommand{\Lt}{L_t}

\pagestyle{fancy}
\fancyhf{}
\lhead{\textsc{Deep Learning} --- IIIT Dharwad}
\rhead{Recurrent Neural Networks: Introduction, BPTT, and Vanishing Gradients}
\rfoot{Page \thepage}
\lfoot{Prof.\ S R M Prasanna \& Dr.\ Sunil Saumya}
\renewcommand{\headrulewidth}{0.4pt}

\title{%
  {\Large\bfseries Deep Learning Course}\\[6pt]
  {\LARGE Recurrent Neural Networks: Introduction, BPTT, and Vanishing Gradients}\\[4pt]
  {\normalsize IIIT Dharwad --- M.Tech Programme}
}
\author{Prof.\ S R M Prasanna \& Dr.\ Sunil Saumya}
\date{Sessions 16, 17}

\begin{document}

\maketitle
\tableofcontents
\clearpage

%% ============================================================
%% Session 16: RNN Introduction and Sequence Modelling
%% ============================================================
\part{Session 16: RNN Introduction and Sequence Modelling}

% ============================================================
%  TITLE PAGE
% ============================================================
\begin{titlepage}
  \centering
  \vspace*{2cm}

  {\LARGE\bfseries Deep Learning Course\\[0.4em]
   Indian Institute of Information Technology, Dharwad}

  \vspace{1.5cm}
  \rule{\linewidth}{0.8pt}\\[0.4em]
  {\Huge\bfseries Session 16}\\[0.5em]
  {\Large\bfseries Recurrent Neural Networks:\\[0.3em]
   Motivation and Sequence Modelling}
  \\[0.4em]\rule{\linewidth}{0.8pt}

  \vspace{1.5cm}

  \begin{tabular}{ll}
    \textbf{Session Number:}  & 16 \\[4pt]
    \textbf{Date:}            & January 11--12, 2026 \\[4pt]
    \textbf{Instructors:}     & Prof.\ S.\ R.\ Mahadeva Prasanna \\
                              & Dr.\ Sunil Saumya \\[4pt]
    \textbf{Slide Decks:}     & \texttt{05-DL-RNN} (slides 1--63) \\
                              & \texttt{07-PyTorchDL-TransferLearning} (slides 1--29) \\[4pt]
    \textbf{Duration:}        & $\approx$ 1 hr 37 min \\[4pt]
    \textbf{Topics Covered:}  & Sequence modelling motivation; \\
                              & Limitations of feedforward networks; \\
                              & FFNN $\to$ RNN architectural transition; \\
                              & Variable-length input problem; \\
                              & Order/position blindness problem \\
  \end{tabular}

  \vfill
  {\small These notes are compiled from the lecture transcript and visual
   extractions for archival and study purposes.}
\end{titlepage}

% ============================================================
%  TABLE OF CONTENTS
% ============================================================

\newpage

% ============================================================
\section{Introduction}
\label{sec:introduction}
% ============================================================

Session~16 marks the transition in the deep learning course from spatial
models (convolutional neural networks, CNNs) to \emph{sequential} models.
Having established CNN architectures and transfer learning in the preceding
sessions, the instructors now motivate a fundamentally different class of
architectures: \textbf{Recurrent Neural Networks (RNNs)} and their
variants---Gated Recurrent Units (GRUs) and Long Short-Term Memory networks
(LSTMs).

The central question posed at the outset is: \emph{Why do feedforward neural
networks (FFNNs), despite their universal approximation capability, fail for
sequential data?} The session answers this question through two distinct
lines of argument:

\begin{enumerate}[leftmargin=2em]
  \item \textbf{Fixed window / long-range dependency failure.} A feedforward
        network operating on a sliding context window is structurally incapable
        of capturing dependencies that span more tokens than the window width.
  \item \textbf{Parallel input / order blindness.} When an entire sequence is
        flattened and fed as a single input vector, the network has no
        mechanism to distinguish word order; permuting the input produces the
        same (or nearly the same) output.
\end{enumerate}

The second half of the session, drawing from the
\texttt{07-PyTorchDL-TransferLearning} deck, formalises these problems in
the concrete setting of text classification with one-hot word vectors,
culminating in the realisation that the input dimension $N \times V$ varies
per sentence---an impossibility for a standard fully connected layer.

\begin{keyinsight}
RNNs were proposed alongside CNNs during the era of shallow neural networks,
but their true effectiveness was only demonstrated with the computational
power and training techniques of the deep learning era.  Both CNN and RNN
are \emph{structured} extensions of the basic MLP: CNNs exploit spatial
locality; RNNs exploit temporal/sequential ordering.
\end{keyinsight}

% ============================================================
\section{Background and Prerequisites}
\label{sec:background}
% ============================================================

\subsection{Data Taxonomy: Sequential vs.\ Non-Sequential}
\label{subsec:data_taxonomy}

Before introducing RNNs, the instructors establish a clear dichotomy between
the two data types that have been encountered so far in the course.

\begin{definition}[Non-Sequential Data]
\label{def:nonseq}
Data is \emph{non-sequential} when the order in which features (or rows) are
presented to a model carries no semantic information.  Permuting features
does not change the classification or regression target.
\end{definition}

\begin{definition}[Sequential Data]
\label{def:seq}
Data is \emph{sequential} when the temporal or positional ordering of
observations encodes information that is necessary to recover the target
label or prediction.  Permuting the observations destroys this information.
\end{definition}

\Cref{tab:data_types} gives concrete examples of both categories encountered
in the lecture.

\begin{table}[ht]
  \centering
  \caption{Non-sequential vs.\ sequential data examples discussed in Session~16.}
  \label{tab:data_types}
  \begin{tabularx}{\linewidth}{lXX}
    \toprule
    \textbf{Type} & \textbf{Example} & \textbf{Key Property} \\
    \midrule
    Non-sequential & Student placement data (CGPA, IQ, placement) &
      Shuffling feature columns does not change prediction \\
    Non-sequential & MNIST digit images &
      Fixed spatial grid; no temporal component \\
    Sequential & Stock price time series (open, high, close by date) &
      Shuffling rows destroys temporal patterns \\
    Sequential & Natural language text &
      Word order encodes meaning \\
    Sequential & Speech signal &
      Phoneme sequence encodes words \\
    Sequential & Video frames &
      Frame order encodes actions and events \\
    \bottomrule
  \end{tabularx}
\end{table}

\subsection{One-Hot Encoding Recap}
\label{subsec:onehot}

Throughout the session, words are represented as \emph{one-hot vectors}.
Given a vocabulary $\mathcal{V}$ of size $V$, each word $w_i$ is mapped to a
vector $\mathbf{e}_i \in \{0,1\}^{V}$ with exactly one non-zero component:

\begin{equation}
  \label{eq:onehot}
  [\mathbf{e}_i]_j =
  \begin{cases}
    1 & \text{if } j = i, \\
    0 & \text{otherwise.}
  \end{cases}
\end{equation}

A sentence of $N$ words is then a sequence
$(\mathbf{e}_{w_1}, \mathbf{e}_{w_2}, \ldots, \mathbf{e}_{w_N})$, each
vector of dimension $V$.

\subsection{Standard FFNN Review}
\label{subsec:ffnn_review}

A standard feedforward neural network with input dimension $d_{\text{in}}$,
one hidden layer of width $H$, and output dimension $d_{\text{out}}$ computes:

\begin{align}
  \label{eq:ffnn_hidden}
  \mathbf{h} &= \sigma\!\left(W_1 \mathbf{x} + \mathbf{b}_1\right), \\
  \label{eq:ffnn_output}
  \hat{\mathbf{y}} &= \text{softmax}\!\left(W_2 \mathbf{h} + \mathbf{b}_2\right),
\end{align}

where $W_1 \in \mathbb{R}^{H \times d_{\text{in}}}$,
$W_2 \in \mathbb{R}^{d_{\text{out}} \times H}$, and $\sigma$ is an
element-wise nonlinearity (e.g., ReLU or tanh).  The architecture is
\emph{fixed}: $d_{\text{in}}$ is a compile-time hyperparameter that cannot
change between inference calls.

% ============================================================
\section{Motivation for Sequence Modelling}
\label{sec:motivation}
% ============================================================

\subsection{The Significance of Sequential Order}
\label{subsec:order_significance}

Prof.\ Prasanna opens the technical discussion with a compelling linguistic
argument.  Consider the Hindi words \emph{kama} and \emph{kamala}: they
share identical constituent phonemes (\emph{ka}, \emph{ma}, \emph{la}), yet
their meanings differ entirely because of the \emph{order} in which the
phonemes appear.  A model based purely on the distribution of phoneme
occurrences would find the two words indistinguishable.

\begin{example}
The English word ``book'' illustrates context-dependent meaning driven by
sequence:
\begin{itemize}[nosep]
  \item ``I want to \textbf{book} a flight.'' \quad (\textit{book} = verb,
        to reserve)
  \item ``I want to read a \textbf{book}.'' \quad (\textit{book} = noun,
        a text)
\end{itemize}
The meaning is disambiguated entirely by the neighbouring words in the
sequence.
\end{example}

\begin{example}[title={Ball Trajectory Prediction}]
Given only the \emph{current position} of a ball on a field, all directions
of motion are equally probable.  However, given the sequence of positions
at $t{=}1,2,3,4$, the direction at $t{=}5$ can be predicted with high
confidence.  This illustrates that \emph{past sequence context} provides
the necessary information.
\end{example}

\subsection{Requirements for a Sequence Modelling Architecture}
\label{subsec:requirements}

The instructors enumerate the wishlist for any architecture designed to
handle sequential data (slide V01, \texttt{05-DL-RNN}, slide 9/63):

\begin{enumerate}[label=\textbf{R\arabic*.}, leftmargin=3em]
  \item \label{req:order} \textbf{Preserve order.} The model must process
        tokens in the order they appear; $w_t$ must be processed after
        $w_{t-1}$.
  \item \label{req:varlen} \textbf{Handle variable sequence lengths.} For the
        same class of information, different instances may have different
        sequence lengths.  The architecture must not require a fixed length.
  \item \label{req:params} \textbf{Parameter count independent of sequence
        length.} Parameters should not grow linearly (or worse) with
        sequence length; they must be \emph{shared} across time steps.
  \item \label{req:longterm} \textbf{Model long-range dependencies.}
        Information from many steps in the past must remain accessible when
        generating an output at the current step.
\end{enumerate}

\begin{keyinsight}
Requirements \ref{req:params} and \ref{req:longterm} are in tension: a model
with shared parameters that still retains long-range information requires a
form of \emph{memory} that persists and evolves across time steps.  This is
precisely the role of the hidden state $h_t$ in an RNN.
\end{keyinsight}

\subsection{Applications of Sequence Modelling}
\label{subsec:applications}

\Cref{tab:applications} summarises the sequence modelling tasks discussed in
the lecture.

\begin{table}[ht]
  \centering
  \caption{Representative sequence modelling tasks and their input/output
           structure.}
  \label{tab:applications}
  \begin{tabularx}{\linewidth}{lXX}
    \toprule
    \textbf{Task} & \textbf{Input $X$} & \textbf{Output $Y$} \\
    \midrule
    Speech recognition &
      Sequence of acoustic feature frames &
      Sequence of recognised words \\
    Machine translation &
      Sentence in source language &
      Sentence in target language \\
    Text summarisation &
      Long document ($\sim$200 words) &
      Compressed summary ($\sim$25 words) \\
    Sentiment classification &
      Text review &
      Sentiment label (positive/negative) \\
    Named entity recognition &
      Token sequence &
      Per-token entity label \\
    Video action recognition &
      Sequence of video frames &
      Action label \\
    Stock price forecasting &
      Time-ordered OHLC prices &
      Next-day closing price \\
    \bottomrule
  \end{tabularx}
\end{table}

% ============================================================
\section{Limitations of Feedforward Networks for Sequential Data}
\label{sec:ffnn_limitations}
% ============================================================

\subsection{Limitation 1: Fixed Window Cannot Model Long-Range Dependencies}
\label{subsec:fixed_window}

One naive approach to sequence modelling with FFNNs is to use a \emph{fixed
sliding window} of $k$ preceding tokens as context.  Given the current
position $t$, the network receives
$(w_{t-k}, w_{t-k+1}, \ldots, w_{t-1})$ and predicts $w_t$.

\begin{warningbox}
A fixed window of size $k$ can only capture dependencies within a span of
$k$ tokens.  Any dependency of length $> k$ is structurally invisible to
the model, regardless of how many parameters are trained.
\end{warningbox}

The lecture illustrates this with the following sentence (from Ava Soleimany,
MIT 6.S191, 2020):

\begin{quote}
``France is where I grew up, but I now live in Boston.
 I speak fluent \underline{\hspace{1.5cm}}.''
\end{quote}

The blank depends on the word ``France'' which appears many tokens earlier.
A short fixed window sees only ``\ldots live in Boston. I speak fluent'' and
might incorrectly predict ``English'' rather than ``French.''

\vspace{0.5em}
The fundamental limitation can be expressed as:

\begin{equation}
  \label{eq:fixed_window}
  P(w_t \mid w_{t-1}, w_{t-2}, \ldots, w_1)
  \;\neq\;
  P(w_t \mid w_{t-1}, w_{t-2}, \ldots, w_{t-k})
  \quad \text{when the dependency length} > k.
\end{equation}

\subsection{Limitation 2: Variable-Length Input Dimension}
\label{subsec:variable_length}

Consider a corpus $\mathcal{C}$ of sentences, each sentence
$s_i = (w_1^{(i)}, w_2^{(i)}, \ldots, w_{N_i}^{(i)})$ having $N_i$ words,
where $N_i$ varies across sentences.  Using one-hot encoding with vocabulary
size $V$, the flattened input for sentence $s_i$ has dimension $N_i \times V$.

\Cref{tab:varlen} reproduces the slide from the
\texttt{07-PyTorchDL-TransferLearning} deck (slide 7/29).

\begin{table}[ht]
  \centering
  \caption{Text input dimension as a function of sentence length.
           $V = $ vocabulary size.}
  \label{tab:varlen}
  \begin{tabular}{lcl}
    \toprule
    \textbf{Input Sentence} & \textbf{\# Words} ($N$) &
    \textbf{Network Input Size} \\
    \midrule
    We study ML               & 3 & $3 \times V$ \\
    We are in LLM class       & 5 & $5 \times V$ \\
    Today we study neural networks & 5 & $5 \times V$ \\
    Deep learning is powerful & 4 & $4 \times V$ \\
    \bottomrule
  \end{tabular}
\end{table}

An ANN requires a \emph{fixed} input dimension $d_{\text{in}}$.  Thus the
input dimension must be set to $d_{\text{in}} = N_{\max} \times V$, where
$N_{\max}$ is the length of the longest sentence in the corpus.  Shorter
sentences are \emph{zero-padded} to length $N_{\max}$.

\subsubsection*{Numerical Scale of the Problem}

Suppose $V = 1000$ unique words and $N_{\max} = 100$ words.  Then:

\begin{align}
  d_{\text{in}} &= N_{\max} \times V = 100 \times 1000 = 100{,}000, \\
  \#\text{weights (input}\to\text{hidden)} &= d_{\text{in}} \times H
    = 100{,}000 \times H.
\end{align}

For even a modest hidden layer of $H = 3$, this yields $3 \times 10^5$
weights---applied \emph{identically} to a sentence with only 5 words (which
contributes a mere $5 \times 1000 = 5{,}000$ non-zero inputs, with the
remaining $95{,}000$ inputs being zero padding).

\begin{keyinsight}
Zero-padding wastes computation proportionally to $(N_{\max} - N_i)/N_{\max}$
for each sentence $s_i$.  In the worst case (a one-word sentence against a
100-word maximum) 99\% of the input is padding noise, yet 99\% of the
weight multiplications are still performed.
\end{keyinsight}

Furthermore, at test time, if a sentence arrives with $N > N_{\max}$ words,
the network has no architectural mechanism to handle it---a hard failure.

\subsection{Limitation 3: Loss of Sequence Order Information}
\label{subsec:order_loss}

Even if the variable-length problem is ``solved'' by zero-padding, the FFNN
still suffers from order blindness.  When the entire sentence is flattened
into a single input vector and fed to the network in one forward pass, the
network has no way to distinguish \emph{which word came first}.

\begin{example}[title={Order Blindness in ANN}]
Consider the sentence ``We are in LLM class'' with vocabulary
$\mathcal{V} = \{\text{We, are, in, LLM, class}\}$.  The network input is
a $25$-dimensional vector (5 words $\times$ 5-dimensional one-hot each),
shown schematically below:

\vspace{0.3em}
\begin{center}
\begin{tikzpicture}[
    word/.style={draw,rounded corners,fill=green!10,minimum width=1.5cm,
                 minimum height=0.6cm,font=\small},
    arr/.style={-{Stealth},thick}
  ]
  \node[word] (w1) at (0,0)   {We};
  \node[word] (w2) at (1.8,0) {are};
  \node[word] (w3) at (3.6,0) {in};
  \node[word] (w4) at (5.4,0) {LLM};
  \node[word] (w5) at (7.2,0) {class};

  \node[draw,fill=yellow!20,minimum width=3cm,minimum height=0.7cm,
        font=\small] (flat) at (3.6,-1.5) {Flattened: 25-dim vector};

  \draw[arr] (w1.south) -- ++(0,-0.5) -| (flat.north west);
  \draw[arr] (w2.south) -- ++(0,-0.5) -| (flat.north);
  \draw[arr] (w3.south) -- (flat.north);
  \draw[arr] (w4.south) -- ++(0,-0.5) -| (flat.north east);
  \draw[arr] (w5.south) -- ++(0,-0.5) -| (flat.east |- flat.north);

  \node[draw,fill=orange!15,minimum width=1.4cm,minimum height=0.6cm,
        font=\small] (h) at (3.6,-3.0) {Hidden: 3};
  \node[draw,fill=red!15,minimum width=1.2cm,minimum height=0.6cm,
        font=\small] (o) at (3.6,-4.2) {Output: 1};

  \draw[arr] (flat.south) -- (h.north);
  \draw[arr] (h.south) -- (o.north);

  \node[font=\footnotesize,text=gray] at (7.5,-1.5)
    {All words fed simultaneously};
  \node[font=\footnotesize,text=gray] at (7.5,-2.0)
    {$\Rightarrow$ order is invisible};
\end{tikzpicture}
\end{center}
\vspace{0.3em}
The sentence ``class LLM in are We'' produces an \emph{identical} input
vector (up to reordering within the flattened representation), so the
network's output is the same regardless of word order.
\end{example}

\begin{remark}
\label{rem:order_loss}
This is a fundamental architectural limitation, not a data or training
deficiency.  No amount of additional data or regularisation can teach an
ANN to distinguish word order when the order information is discarded before
the first matrix multiplication.
\end{remark}

\subsection{Summary of FFNN Limitations for Sequential Data}
\label{subsec:ffnn_summary}

\begin{table}[ht]
  \centering
  \caption{Summary of feedforward network limitations for sequential data.}
  \label{tab:ffnn_limits}
  \begin{tabularx}{\linewidth}{lXX}
    \toprule
    \textbf{Limitation} & \textbf{Description} & \textbf{Consequence} \\
    \midrule
    Fixed input size &
      FFNN requires all inputs to have identical dimension &
      Cannot process variable-length sequences without padding \\
    Long-range dependency &
      Fixed window of size $k$ misses dependencies $> k$ &
      Important contextual cues (e.g., ``France $\to$ French'') are lost \\
    Order blindness &
      Parallel input processing ignores token position &
      Permutations of the same words produce the same output \\
    Parameter explosion &
      Zero-padding inflates input dimension to $N_{\max} \times V$ &
      Computation/memory scale with worst-case sequence length \\
    Poor generalisation &
      Model trained on max length $N_{\max}$ fails on $N > N_{\max}$ &
      Hard failure at test time on unseen longer sequences \\
    No parameter sharing &
      Each position has independent weight matrix in a position-wise MLP &
      The same word at different positions is processed differently \\
    \bottomrule
  \end{tabularx}
\end{table}

% ============================================================
\section{From FFNN to RNN: The Architectural Transition}
\label{sec:ffnn_to_rnn}
% ============================================================

\subsection{Conceptual Transition}
\label{subsec:conceptual_transition}

The core idea behind the FFNN $\to$ RNN transition is to \emph{unfold} the
network across time and then introduce \emph{weight sharing} across positions,
together with a \emph{recurrent connection} that carries information from
the previous time step.

Prof.\ Prasanna explains the progression in three steps:

\begin{enumerate}[leftmargin=2em]
  \item \textbf{Multiple FFNNs side by side.} Place one FFNN per token
        position, processing each token independently.  This is conceptually
        correct for position-by-position processing but does not share
        parameters and grows with sequence length.
  \item \textbf{Collapse to simplified notation.} Represent each FFNN
        (input layer $\to$ hidden layers $\to$ output layer) as a single
        compact symbol $x \to [h] \to y$.
  \item \textbf{Add the recurrent connection.} Instead of replicating the
        compact symbol horizontally (one per position), \emph{rotate it} and
        add a self-loop (arrow $C$) on the hidden node.  This loop carries the
        hidden state $h_{t-1}$ forward to time $t$, sharing the same weights
        across all time steps.
\end{enumerate}

\subsection{FFNN vs.\ RNN Architecture Diagram}
\label{subsec:arch_diagram}

\Cref{fig:ffnn_vs_rnn} reproduces the key comparison from slide V03
(\texttt{05-DL-RNN}, slide 15/63).

\begin{figure}[ht]
  \centering
  \begin{tikzpicture}[
      node distance=1.6cm,
      layer/.style={draw,rounded corners,fill=blue!10,minimum width=1.2cm,
                    minimum height=0.8cm,font=\small\bfseries},
      hlayer/.style={draw,rounded corners,fill=yellow!20,minimum width=1.2cm,
                     minimum height=0.8cm,font=\small\bfseries},
      olayer/.style={draw,rounded corners,fill=red!15,minimum width=1.2cm,
                     minimum height=0.8cm,font=\small\bfseries},
      arr/.style={-{Stealth[length=6pt]},thick},
      label/.style={font=\small\itshape}
    ]

    %% ─── FFNN (top) ──────────────────────────────────────────────────
    \node[font=\bfseries\small] at (-3.5,3.6) {Feedforward Neural Network};
    \node[layer]  (fx) at (-3.5,2.6) {$x$};
    \node[hlayer] (fh) at (-3.5,1.3) {$h$};
    \node[olayer] (fy) at (-3.5,0.0) {$\hat{y}$};

    \draw[arr] (fx) -- (fh) node[midway,right,label]{$W_1$};
    \draw[arr] (fh) -- (fy) node[midway,right,label]{$W_2$};

    \node[font=\footnotesize,text=gray] at (-3.5,-0.7)
      {(simplified: one block per layer)};

    %% ─── RNN (bottom) ────────────────────────────────────────────────
    \node[font=\bfseries\small] at (2.5,3.6) {Recurrent Neural Network};

    % Unrolled version
    \foreach \t/\lbl in {0/t{-}1,2.5/t,5.0/t{+}1}{
      \node[layer]  (rx\t) at (\t+0.0, 2.6) {$x_{\lbl}$};
      \node[hlayer] (rh\t) at (\t+0.0, 1.3) {$h_{\lbl}$};
      \node[olayer] (ry\t) at (\t+0.0, 0.0) {$y_{\lbl}$};

      \draw[arr] (rx\t) -- (rh\t);
      \draw[arr] (rh\t) -- (ry\t);
    }

    % Recurrent connections
    \draw[arr] (rh0) -- (rh2.5) node[midway,above,label]{$\Whh$};
    \draw[arr] (rh2.5) -- (rh5.0) node[midway,above,label]{$\Whh$};

    % Ellipses
    \node at (-1.0,1.3) {$\cdots$};
    \node at (6.6,1.3)  {$\cdots$};

    %% ─── Compact RNN symbol ──────────────────────────────────────────
    \node[font=\bfseries\small] at (10.0,3.6) {Compact form};
    \node[layer]  (cx) at (10.0,2.6) {$x_t$};
    \node[hlayer] (ch) at (10.0,1.3) {$h_t$};
    \node[olayer] (cy) at (10.0,0.0) {$y_t$};

    \draw[arr] (cx) -- (ch) node[midway,right,label]{$B$};
    \draw[arr] (ch) -- (cy) node[midway,right,label]{$A$};

    % Self-loop
    \draw[arr] (ch.east) to[out=0,in=0,looseness=3]
               node[right,label]{$C$ ($\Whh$)} (ch.east |- ch.south)
               to[out=180,in=270] (ch.south);

    % Labels for connections
    \node[font=\footnotesize,text=blue!70!black] at (10.0,-0.8)
      {$A$: output conn.\ ($h \to y$)};
    \node[font=\footnotesize,text=blue!70!black] at (10.0,-1.2)
      {$B$: input conn.\ ($x \to h$)};
    \node[font=\footnotesize,text=blue!70!black] at (10.0,-1.6)
      {$C$: recurrent conn.\ ($h \to h$)};

  \end{tikzpicture}
  \caption{Comparison of FFNN (left) and RNN (right).  The RNN unrolled over
           time (centre) shares weights $\Whh$, $\Wxh$, $\Why$ across all
           time steps.  The compact form (right) uses a self-loop on the
           hidden node to denote the recurrence.}
  \label{fig:ffnn_vs_rnn}
\end{figure}

\subsection{RNN Forward Pass Equations}
\label{subsec:forward_pass}

The RNN computes, at each time step $t$:

\begin{align}
  \label{eq:rnn_hidden}
  \ht &= \tanh\!\left(\Wxh\, \xt + \Whh\, h_{t-1} + \mathbf{b}_h\right), \\
  \label{eq:rnn_output}
  \yt &= \text{softmax}\!\left(\Why\, \ht + \mathbf{b}_y\right),
\end{align}

where:
\begin{itemize}[nosep]
  \item $\xt \in \mathbb{R}^{d_x}$ is the input at time $t$,
  \item $h_{t-1} \in \mathbb{R}^{d_h}$ is the hidden state from the previous
        time step,
  \item $\Wxh \in \mathbb{R}^{d_h \times d_x}$ maps input to hidden,
  \item $\Whh \in \mathbb{R}^{d_h \times d_h}$ maps hidden to hidden
        (the recurrent weight),
  \item $\Why \in \mathbb{R}^{d_y \times d_h}$ maps hidden to output.
\end{itemize}

\begin{keyinsight}
The three weight matrices $\Wxh$, $\Whh$, $\Why$ are \textbf{shared across
all time steps}.  This is the mechanism by which the number of parameters
becomes \emph{independent} of sequence length, satisfying requirement
\ref{req:params} from \cref{subsec:requirements}.
\end{keyinsight}

\subsection{How the RNN Addresses Each FFNN Limitation}

\begin{table}[ht]
  \centering
  \caption{How RNN architecture addresses the limitations of FFNNs.}
  \label{tab:rnn_solutions}
  \begin{tabularx}{\linewidth}{lX}
    \toprule
    \textbf{FFNN Limitation} & \textbf{RNN Solution} \\
    \midrule
    Fixed input size &
      Each token $x_t$ is fed one at a time; no need to specify sequence
      length at network-design time \\
    Long-range dependency &
      $h_t$ accumulates information from all previous steps; theoretically
      unlimited context window (bounded in practice for vanilla RNN) \\
    Order blindness &
      Tokens are processed sequentially; $h_t$ at step $t$ encodes the
      entire history $x_1, \ldots, x_t$ in a fixed-size summary \\
    Parameter explosion &
      $\Wxh, \Whh, \Why$ are shared; parameter count is $O(d_h^2 + d_h d_x)$,
      independent of $N$ \\
    Poor generalisation &
      The same shared weights can process sequences of any length at test time \\
    No parameter sharing &
      Weight sharing is the \emph{defining feature} of the recurrent connection \\
    \bottomrule
  \end{tabularx}
\end{table}

% ============================================================
\section{Detailed Analysis: Text Input and FFNN Problems}
\label{sec:text_input}
% ============================================================

\subsection{Text Representation Pipeline}
\label{subsec:text_pipeline}

Dr.\ Saumya's portion of the session works through a concrete text example
to make the FFNN limitations tangible.  The pipeline for representing a
sentence as input to a neural network is:

\begin{enumerate}[leftmargin=2em]
  \item \textbf{Tokenise} the sentence into individual words.
  \item \textbf{Build a vocabulary} $\mathcal{V}$ of all unique tokens.
  \item \textbf{Assign integer indices} to each token.
  \item \textbf{One-hot encode} each token using \cref{eq:onehot}.
  \item \textbf{Flatten} all one-hot vectors into a single input vector of
        dimension $N \times V$.
\end{enumerate}

\Cref{alg:text_to_input} formalises this pipeline.

\begin{algorithm}[ht]
  \caption{Text-to-Input Vector for an ANN}
  \label{alg:text_to_input}
  \SetAlgoLined
  \KwIn{Sentence $s = (w_1, w_2, \ldots, w_N)$; vocabulary $\mathcal{V}$,
        $|\mathcal{V}| = V$; max length $N_{\max}$}
  \KwOut{Flattened input vector $\mathbf{x} \in \mathbb{R}^{N_{\max} \times V}$}

  \BlankLine
  Initialise $\mathbf{x} \leftarrow \mathbf{0}_{N_{\max} \times V}$
  \tcp*[r]{zero padding by default}

  \For{$t \leftarrow 1$ \KwTo $N$}{
    $i \leftarrow \text{index of } w_t \text{ in } \mathcal{V}$\;
    $\mathbf{e}_t \leftarrow \mathbf{0}_V$\;
    $[\mathbf{e}_t]_i \leftarrow 1$
    \tcp*[r]{one-hot vector}
    $\mathbf{x}[(t-1)V : tV] \leftarrow \mathbf{e}_t$
    \tcp*[r]{insert into flattened vector}
  }
  \tcp*[r]{positions $N+1$ to $N_{\max}$ remain zero-padded}
  \Return $\mathbf{x}$\;
\end{algorithm}

\subsection{Why Parallel Feeding Destroys Order}
\label{subsec:parallel_destroys_order}

When $\mathbf{x}$ is constructed as above and fed to a standard ANN, the
hidden layer computes:

\begin{equation}
  \label{eq:ann_hidden}
  \mathbf{h} = \sigma(W_1 \mathbf{x} + \mathbf{b}_1),
\end{equation}

where $W_1 \in \mathbb{R}^{H \times (N_{\max} V)}$.  Each column of $W_1$
is associated with a specific position-word combination, not with a word
independent of its position.  Crucially, all components of $\mathbf{x}$
enter the summation simultaneously---there is no mechanism by which the
computation at position $t$ can condition on the result of computation at
position $t-1$.

\begin{remark}
This is fundamentally different from the way humans process language.  When
reading ``We are in LLM class'', the human brain processes each word
sequentially, updating an internal representation after each word.  The ANN
instead ``sees'' all five words at once, without any notion of what came
before what.
\end{remark}

\subsection{The Stochastic Model Predecessor: Hidden Markov Models}
\label{subsec:hmm}

Prof.\ Prasanna notes that before deep learning models, \emph{Hidden Markov
Models (HMMs)} were the dominant approach to sequence modelling (e.g., in
speech recognition and NLP).  HMMs encode sequence knowledge probabilistically
via state-transition probabilities:

\begin{equation}
  \label{eq:hmm}
  P(w_1, w_2, \ldots, w_T) = \prod_{t=1}^{T} P(w_t \mid w_{t-1}),
\end{equation}

(under the first-order Markov assumption).  The success of HMMs demonstrated
that sequence knowledge could be modelled formally, inspiring the development
of recurrent neural architectures that could represent richer, learned
sequence dynamics.

% ============================================================
\section{RNN Architecture: Formal Development}
\label{sec:rnn_formal}
% ============================================================

\subsection{Unrolled Computation Graph}
\label{subsec:unrolled}

The RNN processes a sequence $(x_1, x_2, \ldots, x_T)$ by unrolling the
recurrent computation across $T$ time steps.  \Cref{fig:rnn_unrolled}
shows the unrolled computation graph.

\begin{figure}[ht]
  \centering
  \begin{tikzpicture}[
      node distance=2.2cm,
      block/.style={draw,rounded corners,fill=yellow!20,minimum width=1.0cm,
                    minimum height=0.8cm,font=\small},
      input/.style={draw,rounded corners,fill=blue!12,minimum width=0.9cm,
                    minimum height=0.7cm,font=\small},
      output/.style={draw,rounded corners,fill=red!12,minimum width=0.9cm,
                     minimum height=0.7cm,font=\small},
      arr/.style={-{Stealth[length=5pt]},thick},
      lab/.style={font=\scriptsize,midway,above}
    ]

    \node[input]  (x0) at (0.0,0) {$x_1$};
    \node[input]  (x1) at (2.2,0) {$x_2$};
    \node[input]  (x2) at (4.4,0) {$x_3$};
    \node[input]  (x3) at (6.6,0) {$\cdots$};
    \node[input]  (x4) at (8.8,0) {$x_T$};

    \node[block]  (h0) at (0.0,1.6) {$h_1$};
    \node[block]  (h1) at (2.2,1.6) {$h_2$};
    \node[block]  (h2) at (4.4,1.6) {$h_3$};
    \node         (hd) at (6.6,1.6) {$\cdots$};
    \node[block]  (h4) at (8.8,1.6) {$h_T$};

    \node[output] (y0) at (0.0,3.2) {$y_1$};
    \node[output] (y1) at (2.2,3.2) {$y_2$};
    \node[output] (y2) at (4.4,3.2) {$y_3$};
    \node[output] (y4) at (8.8,3.2) {$y_T$};

    % Input to hidden
    \draw[arr] (x0) -- (h0) node[lab,right]{$\Wxh$};
    \draw[arr] (x1) -- (h1) node[lab,right]{$\Wxh$};
    \draw[arr] (x2) -- (h2) node[lab,right]{$\Wxh$};
    \draw[arr] (x4) -- (h4) node[lab,right]{$\Wxh$};

    % Hidden to output
    \draw[arr] (h0) -- (y0) node[lab,right]{$\Why$};
    \draw[arr] (h1) -- (y1) node[lab,right]{$\Why$};
    \draw[arr] (h2) -- (y2) node[lab,right]{$\Why$};
    \draw[arr] (h4) -- (y4) node[lab,right]{$\Why$};

    % Recurrent connections
    \draw[arr] (h0) -- (h1) node[font=\scriptsize,midway,above]{$\Whh$};
    \draw[arr] (h1) -- (h2) node[font=\scriptsize,midway,above]{$\Whh$};
    \draw[arr] (h2) -- (hd) node[font=\scriptsize,midway,above]{$\Whh$};
    \draw[arr] (hd) -- (h4) node[font=\scriptsize,midway,above]{$\Whh$};

    % Initial hidden state
    \node[font=\small] (h_init) at (-1.6,1.6) {$h_0{=}\mathbf{0}$};
    \draw[arr] (h_init) -- (h0) node[font=\scriptsize,midway,above]{$\Whh$};

    % Annotation
    \draw[decorate,decoration={brace,amplitude=6pt,mirror},
          draw=blue!60,thick]
      (0.0,-0.6) -- (8.8,-0.6)
      node[midway,below=8pt,font=\footnotesize,text=blue!60]
        {same $\Wxh$, $\Whh$, $\Why$ shared across all $T$ steps};

  \end{tikzpicture}
  \caption{Unrolled RNN computation graph.  The initial hidden state $h_0$
           is typically set to zero.  At each step $t$, the hidden state
           $h_t$ is computed from the current input $x_t$ and the previous
           hidden state $h_{t-1}$ using the \emph{same} weight matrices
           $\Wxh$ and $\Whh$.}
  \label{fig:rnn_unrolled}
\end{figure}

\subsection{Hidden State as Sequence Memory}
\label{subsec:hidden_state_memory}

The hidden state $h_t$ is the core mechanism by which an RNN remembers
past context.  Substituting \cref{eq:rnn_hidden} recursively:

\begin{align}
  h_1 &= \tanh(\Wxh x_1 + \Whh h_0 + \mathbf{b}_h), \\
  h_2 &= \tanh(\Wxh x_2 + \Whh h_1 + \mathbf{b}_h), \\
      &\;\vdots \notag \\
  h_t &= \tanh\!\Bigl(\Wxh x_t + \Whh \tanh\!\bigl(\Wxh x_{t-1} + \cdots
         \bigr) + \mathbf{b}_h\Bigr).
\end{align}

Thus $h_t$ is a \emph{nonlinear function of all past inputs}
$x_1, \ldots, x_t$.  This is the sense in which the RNN has ``memory.''

\begin{keyinsight}
The memory of an RNN is not an explicit data structure---it is \emph{encoded
in the activations of the hidden state}.  The hidden state $h_t$ is a
$d_h$-dimensional summary of the entire history $x_1, \ldots, x_t$.  The
quality of this summary---and how far back in time it can reach---is
determined by $d_h$ and the spectral properties of $\Whh$.
\end{keyinsight}

\subsection{Sequence Modelling Tasks and RNN Variants}
\label{subsec:rnn_variants}

As introduced by Prof.\ Prasanna, the RNN family includes several variants
that address specific shortcomings of the vanilla architecture.
\Cref{tab:rnn_variants} provides an overview.

\begin{table}[ht]
  \centering
  \caption{Overview of the RNN family.}
  \label{tab:rnn_variants}
  \begin{tabularx}{\linewidth}{lXX}
    \toprule
    \textbf{Model} & \textbf{Key Feature} & \textbf{Primary Use Case} \\
    \midrule
    Vanilla RNN &
      Single recurrent connection; $h_t = \tanh(\Wxh x_t + \Whh h_{t-1})$ &
      Short-term sequence modelling \\
    GRU (Gated Recurrent Unit) &
      Update gate and reset gate control information flow &
      Medium-range dependencies; fewer params than LSTM \\
    LSTM (Long Short-Term Memory) &
      Separate cell state $c_t$ and hidden state $h_t$; three gates &
      Long-range dependencies; speech, text \\
    Bidirectional RNN &
      Two RNNs: left-to-right and right-to-left; outputs concatenated &
      Tasks requiring full sentence context (e.g., NER) \\
    TDNN (Time Delay Neural Network) &
      Fixed dilation pattern on input; no true recurrence &
      Acoustic modelling in speech \\
    \bottomrule
  \end{tabularx}
\end{table}

\begin{remark}
As discussed in the session, vanilla RNNs still struggle with very
long-range dependencies in practice due to the \emph{vanishing gradient
problem}: gradients of the loss with respect to early time steps decay
exponentially as they are backpropagated through many tanh nonlinearities
and the recurrent weight matrix $\Whh$.  LSTM was specifically designed to
address this by introducing a cell state that can carry gradients through
long sequences with minimal decay.
\end{remark}

% ============================================================
\section{Sequence Types and Input/Output Configurations}
\label{sec:io_configs}
% ============================================================

RNNs are flexible in their input/output structure.  The lecture implicitly
introduces several configurations through the application examples.

\begin{table}[ht]
  \centering
  \caption{RNN input/output configurations.}
  \label{tab:io_configs}
  \begin{tabularx}{\linewidth}{lllX}
    \toprule
    \textbf{Configuration} & \textbf{Input} & \textbf{Output} & \textbf{Example} \\
    \midrule
    One-to-one   & Single $x$      & Single $y$       & Standard classification \\
    One-to-many  & Single $x$      & Sequence $y_{1:T}$ & Image captioning \\
    Many-to-one  & Sequence $x_{1:T}$ & Single $y$    & Sentiment classification \\
    Many-to-many (sync) & Seq $x_{1:T}$ & Seq $y_{1:T}$ & POS tagging, NER \\
    Many-to-many (async) & Seq $x_{1:T}$ & Seq $y_{1:T'}$ & Machine translation \\
    \bottomrule
  \end{tabularx}
\end{table}

% ============================================================
\section{Summary and Conclusion}
\label{sec:conclusion}
% ============================================================

Session~16 laid the motivational foundation for the entire RNN module of the
course.  The key takeaways are:

\begin{enumerate}[leftmargin=2em,label=\textbf{\arabic*.}]
  \item \textbf{Sequential data is pervasive.} Speech, text, video, and
        financial time series all require models that can exploit the ordering
        of observations.  This is not a niche requirement---it is central to
        the majority of modern AI applications.

  \item \textbf{FFNNs fail for sequential data on three fronts.}
        \begin{itemize}[nosep]
          \item Fixed input size prevents processing variable-length sequences.
          \item Parallel input processing destroys word order.
          \item Zero-padding ``fixes'' the length problem but at enormous
                computational cost and with poor generalisation to unseen lengths.
        \end{itemize}

  \item \textbf{The RNN solves all three problems.}  By processing one token
        at a time with shared weights and a recurrent hidden state, the RNN:
        (a) handles arbitrary sequence lengths, (b) preserves order (since
        $h_t$ depends on $h_{t-1}$), and (c) keeps parameter count
        $O(d_h^2 + d_h d_x)$ independent of sequence length.

  \item \textbf{The hidden state $h_t$ is the key novelty.}  It is a
        $d_h$-dimensional running summary of the sequence up to time $t$,
        updated at every step by the same matrix $\Whh$.

  \item \textbf{Vanilla RNN is a starting point.}  GRU and LSTM extend the
        vanilla RNN with gating mechanisms to address the vanishing gradient
        problem and improve long-range memory.  These will be the subject of
        subsequent sessions.
\end{enumerate}

\begin{keyinsight}
The word ``recurrent'' in \emph{Recurrent Neural Network} refers to the
\emph{recurrence} (repetition/sharing) of the same weight matrix $\Whh$
across every position in the sequence.  This is the architectural realisation
of the principle that the same learned transformation should apply at every
time step---analogous to how convolutional filters are shared spatially in a
CNN.
\end{keyinsight}

\vspace{1.5em}
\noindent\textbf{Preview of Upcoming Sessions.}  The next session will
derive the full forward and backward passes of the vanilla RNN, introduce the
Backpropagation Through Time (BPTT) algorithm, and demonstrate the vanishing
gradient problem quantitatively.  Subsequent sessions will cover GRU, LSTM,
bidirectional RNNs, and ultimately attention mechanisms and the Transformer
architecture.

% ─── Bibliography placeholder ──────────────────────────────────────────────
\vspace{2em}
\noindent\rule{\linewidth}{0.4pt}
\begin{small}
\noindent\textbf{References and Source Materials}\\[3pt]
Slide deck \texttt{05-DL-RNN}, January 11, 2026.
Sunil Saumya \& S.\ R.\ Mahadeva Prasanna, IIITDwd.\\[2pt]
Slide deck \texttt{07-PyTorchDL-TransferLearning}, January 12, 2026.
Sunil Saumya \& S.\ R.\ Mahadeva Prasanna, IIITDwd.\\[2pt]
Ava Soleimany, ``Deep Sequence Modeling,'' MIT 6.S191 (2020):
Recurrent Neural Networks.\\[2pt]
Lecture transcript: Vimeo \texttt{https://vimeo.com/1153669483},
Session~16, IIITDwd Deep Learning Course, January 2026.
\end{small}

\clearpage

%% ============================================================
%% Session 17: Plain RNN, Backpropagation Through Time, and Vanishing Gradients
%% ============================================================
\part{Session 17: Plain RNN, Backpropagation Through Time, and Vanishing Gradients}

%% ══════════════════════════════════════════════

%% ──────────────────────────────────────────────
%% Title page
%% ──────────────────────────────────────────────
\begin{titlepage}
  \centering
  \vspace*{2cm}
  {\LARGE\bfseries Deep Learning Course\\[0.4em]
   Session 17 Lecture Notes\par}
  \vspace{1.5cm}
  {\large\bfseries Recurrent Neural Networks:\\
  Parameter Sharing, Computational Graphs,\\
  Backpropagation Through Time, and the\\
  Vanishing Gradient Problem\par}
  \vspace{1.5cm}
  \begin{tabular}{ll}
    \textbf{Session Number:} & 17 \\[0.3em]
    \textbf{Date:}           & January 13, 2026 \\[0.3em]
    \textbf{Duration:}       & $\approx$ 1 hour 55 minutes \\[0.3em]
    \textbf{Instructors:}    & Prof.\ S R M Prasanna \& Dr.\ Sunil Saumya \\[0.3em]
    \textbf{Institution:}    & IIIT Dharwad \\[0.3em]
    \textbf{Slide Deck:}     & \texttt{05-DL-RNN} (slides 17--43) \\[0.3em]
    \textbf{Video:}          & \url{https://vimeo.com/1155843264} \\
  \end{tabular}
  \vspace{2cm}
  \hrule
  \vspace{0.5cm}
  {\itshape Topics covered: RNN parameter-sharing rationale; RNN computational
  graph across time (BPTT); vanishing gradient problem; RNN types and
  applications; motivation for LSTM/GRU.\par}
  \vspace{0.5cm}
  \hrule
  \vfill
  {\small These notes are compiled from lecture transcripts and slide visual
  extractions.  They are intended as a study aid and do not replace the
  official course material.}
\end{titlepage}

%% ──────────────────────────────────────────────
%% Table of contents
%% ──────────────────────────────────────────────

\newpage

%% ══════════════════════════════════════════════
\section{Introduction}
\label{sec:introduction}
%% ══════════════════════════════════════════════

Session~17 of the deep learning course continues the study of sequence
modelling that was introduced in Session~16.  The session opens with a recap of
why standard feedforward (fully connected) networks fail on sequential data and
then motivates the key innovation of Recurrent Neural Networks (RNNs):
\emph{parameter sharing across time steps}.

The main teaching arc of the session is:

\begin{enumerate}[label=\arabic*.]
  \item \textbf{Recap of MLP limitations} for sequence data --
        fixed input size, loss of order information, and inefficient
        computation.
  \item \textbf{Parameter sharing as the core RNN idea} -- why the same weight
        matrices are reused at every time step, and what benefits this brings.
  \item \textbf{Folded and unfolded RNN representations} -- how the compact
        loop notation expands into a feed-forward computation graph in time.
  \item \textbf{RNN forward equations} -- the three weight matrices
        $\Wxh$, $\Whh$, and $\Why$ and how they combine to produce hidden
        states and outputs.
  \item \textbf{RNN Loss and Backpropagation Through Time (BPTT)} -- how the
        total loss is accumulated across time steps and how gradients flow
        backward through the unrolled graph.
  \item \textbf{Vanishing gradient problem in RNNs} -- why both the vertical
        (depth) and horizontal (time) gradient paths degrade, and why this
        makes long-range dependency learning nearly impossible with a plain RNN.
  \item \textbf{Motivation for LSTM and GRU} -- gated architectures that
        provide explicit gradient highways to solve the vanishing gradient
        problem.
\end{enumerate}

\begin{keyinsight}[title=Session Roadmap]
Session~17 covers the plain (vanilla) RNN only.  LSTM is introduced in
Sessions~18--20 and GRU in Sessions~21--22.  Understanding the failures of the
plain RNN is essential for appreciating why gated architectures were invented.
\end{keyinsight}

%% ══════════════════════════════════════════════
\section{Background and Prerequisites}
\label{sec:background}
%% ══════════════════════════════════════════════

\subsection{Sequential vs.\ Non-Sequential Data}
\label{subsec:seq_vs_nonseq}

A \emph{non-sequential} dataset (e.g., a table of patient measurements) treats
each sample as independent.  The order in which rows appear carries no meaning.
A \emph{sequential} dataset is one where order is semantically significant:
swapping the tokens of a sentence, the frames of a video, or the time steps of
a speech signal destroys the meaning.

\begin{definition}[Sequential Data]
\label{def:sequential}
A dataset $\{x^{(1)}, x^{(2)}, \ldots, x^{(N)}\}$ is \emph{sequential} if the
ordering of elements within each sample $x^{(i)} = (x^{(i)}_0, x^{(i)}_1,
\ldots, x^{(i)}_{T_i})$ carries information that must be preserved during
modelling.  The length $T_i$ may differ across samples.
\end{definition}

Examples of sequential data encountered in the course include:

\begin{itemize}[noitemsep]
  \item \textbf{Natural language} -- sentences, paragraphs, documents.
  \item \textbf{Speech} -- sequences of acoustic feature vectors (e.g., MFCCs).
  \item \textbf{Time series} -- stock prices, sensor readings, ECG signals.
  \item \textbf{Video} -- ordered sequences of image frames.
\end{itemize}

\subsection{MLP Limitations for Sequence Modelling}
\label{subsec:mlp_limits}

A multi-layer perceptron (MLP) expects a \emph{fixed-size} input vector.  When
faced with variable-length sequences, a common workaround is \emph{zero
padding}: all sequences are padded to the length of the longest sequence in the
corpus.

\begin{warningbox}[title=Problems with Zero Padding in MLPs]
\begin{enumerate}[label=\arabic*., noitemsep]
  \item \textbf{Computational cost:} The input dimension equals the maximum
        sequence length in the entire training corpus, leading to a massive
        weight matrix even when most inputs are zero.
  \item \textbf{Generalisation failure:} If a test sentence is longer than any
        training sentence, the fixed network cannot handle it.
  \item \textbf{Loss of order information:} All input positions are processed
        simultaneously and symmetrically; the network cannot distinguish which
        token came first.
  \item \textbf{No shared statistical strength:} A pattern (e.g., the word
        ``not'') that occurs at position 3 in one sentence and position 7 in
        another is learned as two completely separate features because each
        position has its own independent weight column.
\end{enumerate}
\end{warningbox}

The instructor summarised the MLP shortcomings as: \emph{fixed input size,
variable-length sequential data, loss of order information, inefficient
computation, and poor generalisation}.

\subsection{RNN Design Goals}
\label{subsec:rnn_goals}

The ideal sequence model should:

\begin{itemize}[noitemsep]
  \item Handle variable-length inputs without modification.
  \item Preserve the order of tokens as information.
  \item Share parameters across positions (time steps) so that a pattern
        appearing at any position is recognised by the same detector.
  \item Maintain a ``memory'' of past inputs when producing each output.
\end{itemize}

These four requirements are precisely what the RNN architecture addresses.

%% ══════════════════════════════════════════════
\section{RNN Architecture and Parameter Sharing}
\label{sec:rnn_arch}
%% ══════════════════════════════════════════════

\subsection{The Parameter-Sharing Principle}
\label{subsec:param_sharing}

\begin{keyinsight}[title=Core Innovation of RNNs]
The key idea in moving from an MLP to an RNN is \textbf{sharing parameters
across different parts of the model} -- specifically, reusing the \emph{same}
weight matrices at every time step.  In an MLP every link has a unique,
independently-learned weight; in an RNN a single set of weight matrices
$\{\Wxh, \Whh, \Why\}$ is applied repeatedly as the network processes each
token in the sequence.
\end{keyinsight}

Why does sharing parameters help?

\begin{enumerate}[label=\arabic*.]
  \item \textbf{Variable-length sequences:} Because the same structure is
        applied at every time step, the network can loop for as many steps as
        the current sequence requires -- 3 steps for a 3-word sentence, 200
        steps for a 200-word document -- without any architectural change.
  \item \textbf{Generalisation across positions:} If the word ``not'' reverses
        sentiment whether it appears at position 2 or position 12, the shared
        weights learn this once and apply it everywhere.
  \item \textbf{Reduced parameter count:} Instead of $T \times d_x \times d_h$
        parameters for $T$ independent position-specific weight matrices, there
        are only $d_x \times d_h$ parameters in $\Wxh$, shared across all $T$
        positions.
  \item \textbf{Statistical strength:} Every occurrence of a pattern in the
        training corpus -- regardless of its position -- contributes gradient
        signal to the same weight update.
\end{enumerate}

\begin{remark}[RNN vs.\ Feedback Network]
\label{rem:feedforward}
Although RNN diagrams show a loop arrow, an RNN is still a \emph{feed-forward}
architecture: signal flows from input to output at each time step.  The loop
represents the fact that the weight state computed at time $t$ is \emph{passed
forward} to time $t+1$; it does not mean signal flows backwards through the
network in the way that a true feedback (recurrent) network would.
\end{remark}

\subsection{Folded and Unfolded RNN Representations}
\label{subsec:folded_unfolded}

\Cref{fig:rnn_compact} shows the \emph{folded} (compact) representation of a
single-layer RNN.  The recurrent cell $h$ represents the entire hidden-layer
structure.  The self-loop on $h$ indicates that the hidden state at time $t$
depends on the hidden state at time $t-1$.

\begin{figure}[ht]
\centering
\begin{tikzpicture}[
    node distance=1.8cm,
    cell/.style={draw, rectangle, rounded corners=4pt, minimum width=1.6cm,
                 minimum height=1.0cm, fill=blue!15, thick},
    io/.style={draw, circle, minimum size=0.85cm, fill=orange!20, thick},
    arr/.style={-{Stealth[length=7pt]}, thick},
    label/.style={font=\small}
  ]
  % Nodes
  \node[io]  (x)  {$x_t$};
  \node[cell, right=of x] (h) {$h_t$};
  \node[io,  right=of h]  (y) {$\hat{y}_t$};

  % Forward arrows
  \draw[arr] (x) -- node[above,label]{$\Wxh$} (h);
  \draw[arr] (h) -- node[above,label]{$\Why$} (y);

  % Self-loop (recurrent connection)
  \draw[arr] (h.north) to[out=120,in=60,looseness=3.5]
             node[above,label]{$\Whh$} (h.north west);

\end{tikzpicture}
\caption{Folded (compact) representation of a vanilla RNN.  The loop on the
hidden cell $h_t$ indicates that the previous hidden state $h_{t-1}$ (carrying
sequence knowledge) is fed back as an additional input at the next time step.}
\label{fig:rnn_compact}
\end{figure}

\Cref{fig:rnn_unfolded} shows the same RNN \emph{unfolded} across time.  Each
vertical slice is identical to the folded network; the horizontal arrows carry
the hidden state forward.  Crucially, the weight matrices $\Wxh$, $\Whh$, and
$\Why$ are \emph{identical} in every slice -- this is the parameter sharing.

\begin{figure}[ht]
\centering
\begin{tikzpicture}[
    node distance=0.9cm and 2.2cm,
    cell/.style={draw, rectangle, rounded corners=4pt,
                 minimum width=1.3cm, minimum height=0.9cm,
                 fill=blue!15, thick},
    io/.style={draw, circle, minimum size=0.75cm,
               fill=orange!20, thick, font=\small},
    loss/.style={draw, diamond, minimum size=0.75cm,
                 fill=red!20, thick, font=\small},
    arr/.style={-{Stealth[length=6pt]}, thick},
    darr/.style={-{Stealth[length=6pt]}, thick, dashed},
    lbl/.style={font=\footnotesize}
  ]
  % Time step 0
  \node[io]  (x0) {$x_0$};
  \node[cell, above=of x0] (h0) {$h_0$};
  \node[io,   above=of h0] (y0) {$\hat{y}_0$};
  \node[loss, above=of y0] (L0) {$L_0$};

  % Time step 1
  \node[io,   right=2.4cm of x0] (x1) {$x_1$};
  \node[cell, above=of x1]       (h1) {$h_1$};
  \node[io,   above=of h1]       (y1) {$\hat{y}_1$};
  \node[loss, above=of y1]       (L1) {$L_1$};

  % Time step 2
  \node[io,   right=2.4cm of x1] (x2) {$x_2$};
  \node[cell, above=of x2]       (h2) {$h_2$};
  \node[io,   above=of h2]       (y2) {$\hat{y}_2$};
  \node[loss, above=of y2]       (L2) {$L_2$};

  % Ellipsis / time T
  \node[io,   right=2.4cm of x2] (xT) {$x_T$};
  \node[cell, above=of xT]       (hT) {$h_T$};
  \node[io,   above=of hT]       (yT) {$\hat{y}_T$};
  \node[loss, above=of yT]       (LT) {$L_T$};

  % Total loss node
  \node[loss, fill=red!40, above=0.6cm of L1, xshift=1.2cm] (L) {$L$};

  % Input -> hidden
  \draw[arr] (x0) -- node[right,lbl]{$\Wxh$} (h0);
  \draw[arr] (x1) -- node[right,lbl]{$\Wxh$} (h1);
  \draw[arr] (x2) -- node[right,lbl]{$\Wxh$} (h2);
  \draw[arr] (xT) -- node[right,lbl]{$\Wxh$} (hT);

  % Hidden -> output
  \draw[arr] (h0) -- node[right,lbl]{$\Why$} (y0);
  \draw[arr] (h1) -- node[right,lbl]{$\Why$} (y1);
  \draw[arr] (h2) -- node[right,lbl]{$\Why$} (y2);
  \draw[arr] (hT) -- node[right,lbl]{$\Why$} (yT);

  % Output -> loss
  \draw[arr] (y0) -- (L0);
  \draw[arr] (y1) -- (L1);
  \draw[arr] (y2) -- (L2);
  \draw[arr] (yT) -- (LT);

  % Loss -> total
  \draw[arr] (L0) -- (L);
  \draw[arr] (L1) -- (L);
  \draw[arr] (L2) -- (L);
  \draw[arr] (LT) -- (L);

  % Recurrent connections h -> h (shared Whh)
  \draw[arr] (h0) -- node[above,lbl]{$\Whh$} (h1);
  \draw[arr] (h1) -- node[above,lbl]{$\Whh$} (h2);
  \draw[darr](h2) -- node[above,lbl]{$\cdots$} (hT);

  % Initial state
  \node[left=0.7cm of h0, font=\small] (h_init) {$h_{-1}$};
  \draw[arr] (h_init) -- (h0);

\end{tikzpicture}
\caption{Unrolled RNN showing the computation graph across time (BPTT diagram).
The same weight matrices $\Wxh$, $\Whh$, $\Why$ appear at \emph{every}
time step, illustrating parameter sharing.  Each time step produces a local
loss $L_t$; the total loss $L = \sum_t L_t$ accumulates these.  Gradients flow
\emph{backward through time} (right to left along the $\Whh$ arrows) as well
as vertically through each slice.}
\label{fig:rnn_unfolded}
\end{figure}

\subsection{Forward Pass Equations}
\label{subsec:forward_pass}

At every time step $t$, the RNN performs two computations:

\begin{align}
  \ht &= \tanh\!\bigl(\Wxh\, \xt + \Whh\, \htmo + b_h\bigr)
  \label{eq:hidden_state} \\[4pt]
  \yt &= \mathrm{softmax}\!\bigl(\Why\, \ht + b_y\bigr)
  \label{eq:output}
\end{align}

where:

\begin{itemize}[noitemsep]
  \item $\xt \in \mathbb{R}^{d_x}$ is the input vector at time $t$ (e.g., a
        one-hot or embedding vector for a word).
  \item $\htmo \in \mathbb{R}^{d_h}$ is the hidden state from the previous
        time step (the ``memory'').
  \item $\Wxh \in \mathbb{R}^{d_h \times d_x}$ maps input to hidden.
  \item $\Whh \in \mathbb{R}^{d_h \times d_h}$ maps previous hidden state to
        current hidden state (the recurrent weight).
  \item $\Why \in \mathbb{R}^{d_y \times d_h}$ maps hidden state to output.
  \item $b_h \in \mathbb{R}^{d_h}$ and $b_y \in \mathbb{R}^{d_y}$ are bias
        vectors.
  \item $\tanh$ is applied element-wise as the hidden-layer activation.
\end{itemize}

\begin{remark}[Comparison with MLP]
\label{rem:mlp_compare}
The MLP computes $h_t = f(W\, x_t)$ -- there is no $\Whh\, \htmo$ term.
The entire content of the RNN innovation is adding this \emph{one term}: the
previous hidden state, weighted by $\Whh$, carries sequence knowledge forward.
\end{remark}

\subsection{Parameter Count Example}
\label{subsec:param_count}

The lecture walked through a concrete example with $d_x = 4$ (vocabulary of 4
words, one-hot encoded), $d_h = 3$ (3 neurons in the hidden/output layer, no
separate hidden layer for simplicity), and $d_y = 3$.  The total number of
trainable parameters is:

\begin{align}
  \text{Parameters} &= \underbrace{d_x \times d_h}_{\Wxh}
                     + \underbrace{d_h \times d_h}_{\Whh}
                     + \underbrace{d_h}_{\text{bias}} \notag \\
                    &= (4 \times 3) + (3 \times 3) + 3 \notag \\
                    &= 12 + 9 + 3 = 24
  \label{eq:param_count}
\end{align}

These 24 parameters are \emph{shared} across all time steps.  If the sequence
has $T$ tokens, the unrolled network uses the same 24 parameters at all $T$
positions -- it does not multiply by $T$.

\begin{example}[title=Variable-Length Handling]
Suppose we have three training sentences of lengths 3, 5, and 4 words.  With
the RNN architecture above:
\begin{itemize}[noitemsep]
  \item Sentence 1 (3 words): the recurrent loop executes 3 times.
  \item Sentence 2 (5 words): the recurrent loop executes 5 times.
  \item Sentence 3 (4 words): the recurrent loop executes 4 times.
\end{itemize}
In every case the parameter count remains 24.  No padding or architectural
modification is needed.  This is the variable-length advantage of parameter
sharing.
\end{example}

%% ══════════════════════════════════════════════
\section{Backpropagation Through Time (BPTT)}
\label{sec:bptt}
%% ══════════════════════════════════════════════

\subsection{Loss Function}
\label{subsec:loss}

For a many-to-many RNN (one output per time step), the loss at time step $t$ is
typically the cross-entropy between the predicted output $\hat{y}_t$ and the
true label $y_t$:

\begin{equation}
  L_t = -\sum_{k} y_t^{(k)} \log \hat{y}_t^{(k)}
  \label{eq:local_loss}
\end{equation}

The \emph{total loss} over the entire sequence is the sum of individual losses:

\begin{equation}
  L = \sum_{t=1}^{T} L_t
  \label{eq:total_loss}
\end{equation}

This matches the visual extraction slide which showed $L = L_0 + L_1 + \cdots
+ L_T$ with all partial losses feeding into a single total-loss node at the
top of the computation graph.

\subsection{Two Gradient Paths}
\label{subsec:two_paths}

The unrolled RNN has two distinct paths along which gradients must travel
during backpropagation:

\begin{enumerate}[label=\arabic*.]
  \item \textbf{Vertical path (conventional backpropagation):} At each time
        step $t$, the error at $\hat{y}_t$ propagates downward through the
        $\Why$ weight, through the hidden layer, and through $\Wxh$ to the
        input.  This is the same gradient path as in any deep MLP.
  \item \textbf{Horizontal path (BPTT):} The error must also propagate
        \emph{backwards across time}, from step $T$ back to step $0$, through
        the chain of $\Whh$ matrices.  This is unique to sequence models.
\end{enumerate}

\begin{keyinsight}[title=What BPTT Computes]
The gradient of the total loss with respect to the shared weight matrix $W$ is
the \emph{sum} of the per-step gradients:
\begin{equation}
  \frac{\partial L}{\partial W} = \sum_{t=1}^{T}
    \frac{\partial L_t}{\partial W}
  \label{eq:bptt_gradient}
\end{equation}
Each term $\partial L_t / \partial W$ itself requires unrolling the gradient
back to time step $0$, making the full gradient computation $O(T)$.
\end{keyinsight}

\subsection{BPTT Algorithm}
\label{subsec:bptt_algo}

\begin{algorithm}[H]
\caption{Backpropagation Through Time (BPTT)}
\label{alg:bptt}
\SetAlgoLined
\DontPrintSemicolon
\KwIn{Input sequence $\{x_0, x_1, \ldots, x_T\}$, true labels
      $\{y_0, y_1, \ldots, y_T\}$, initial hidden state $h_{-1}$,\\
      \hspace{1.2cm}weight matrices $\Wxh$, $\Whh$, $\Why$}
\KwOut{Gradients $\nabla_{\Wxh} L$, $\nabla_{\Whh} L$, $\nabla_{\Why} L$}
\BlankLine
\tcp{Forward pass}
\For{$t \leftarrow 0$ \KwTo $T$}{
  $h_t \leftarrow \tanh\!\bigl(\Wxh\, x_t + \Whh\, h_{t-1} + b_h\bigr)$\;
  $\hat{y}_t \leftarrow \mathrm{softmax}\!\bigl(\Why\, h_t + b_y\bigr)$\;
  $L_t \leftarrow \mathrm{CrossEntropy}(\hat{y}_t,\, y_t)$\;
}
$L \leftarrow \sum_{t=0}^{T} L_t$\;
\BlankLine
\tcp{Initialise accumulated gradients}
$\nabla_{\Wxh} L \leftarrow \mathbf{0}$;\quad
$\nabla_{\Whh} L \leftarrow \mathbf{0}$;\quad
$\nabla_{\Why} L \leftarrow \mathbf{0}$\;
$\delta_h^{(T+1)} \leftarrow \mathbf{0}$\tcp*{no future gradient at last step}
\BlankLine
\tcp{Backward pass through time}
\For{$t \leftarrow T$ \KwTo $0$}{
  $\delta_y^{(t)} \leftarrow \hat{y}_t - y_t$
    \tcp*{gradient of cross-entropy + softmax}
  $\nabla_{\Why} L \leftarrow \nabla_{\Why} L
    + \delta_y^{(t)}\, h_t^{\top}$\;
  $\delta_h^{(t)} \leftarrow \bigl(\Why^{\top}\, \delta_y^{(t)}
    + \Whh^{\top}\, \delta_h^{(t+1)}\bigr)
    \odot \tanh'(h_t)$\;
  $\nabla_{\Wxh} L \leftarrow \nabla_{\Wxh} L
    + \delta_h^{(t)}\, x_t^{\top}$\;
  $\nabla_{\Whh} L \leftarrow \nabla_{\Whh} L
    + \delta_h^{(t)}\, h_{t-1}^{\top}$\;
}
\Return $\nabla_{\Wxh} L$,\; $\nabla_{\Whh} L$,\; $\nabla_{\Why} L$\;
\end{algorithm}

\subsection{Computational Complexity}
\label{subsec:bptt_complexity}

\begin{table}[ht]
\centering
\caption{Comparison of standard backpropagation (MLP) vs.\ BPTT (RNN).}
\label{tab:bp_vs_bptt}
\renewcommand{\arraystretch}{1.35}
\begin{tabular}{@{}lll@{}}
\toprule
\textbf{Property} & \textbf{Standard BP (MLP)} & \textbf{BPTT (RNN)} \\
\midrule
Gradient paths       & Vertical only             & Vertical \emph{and} horizontal \\
Number of weight copies & $L$ (one per layer)    & 1 (shared across $T$ steps) \\
Memory for activations & $O(L \cdot d_h)$        & $O(T \cdot d_h)$ \\
Computation time     & $O(L)$ per sample         & $O(T \cdot L)$ per sample \\
Vanishing gradient   & Through depth $L$         & Through depth $L$ \emph{and} time $T$ \\
Parallelism          & Limited by depth          & None across time steps \\
\bottomrule
\end{tabular}
\end{table}

%% ══════════════════════════════════════════════
\section{The Vanishing Gradient Problem in RNNs}
\label{sec:vanishing}
%% ══════════════════════════════════════════════

\subsection{Mathematical Origin}
\label{subsec:vanishing_math}

To update $\Whh$ correctly, the gradient must flow from a loss at time $t$
back to hidden state $h_k$ for $k \ll t$.  This requires computing:

\begin{equation}
  \frac{\partial h_t}{\partial h_k}
  = \prod_{i=k+1}^{t} \frac{\partial h_i}{\partial h_{i-1}}
  = \prod_{i=k+1}^{t} \Whh^{\top} \cdot
    \mathrm{diag}\!\bigl(\tanh'(h_{i-1})\bigr)
  \label{eq:gradient_product}
\end{equation}

This is a product of $(t - k)$ matrices, each equal to $\Whh^{\top}$ scaled by
the Jacobian of the $\tanh$ activation.  Two facts drive the vanishing gradient
problem:

\begin{enumerate}[label=\arabic*.]
  \item $\tanh'(x) = 1 - \tanh^2(x) \leq 1$, with equality only at $x = 0$.
        For saturated activations the derivative is nearly zero.
  \item If the largest singular value of $\Whh$ satisfies $\sigma_1(\Whh) < 1$,
        then $\|\Whh^{\top}\|^{(t-k)} \to 0$ exponentially as $(t-k)$ grows.
\end{enumerate}

Combining both effects:

\begin{equation}
  \left\|\frac{\partial h_t}{\partial h_k}\right\| \leq
  \left(\sigma_1(\Whh)\cdot \max_x |\tanh'(x)|\right)^{t-k}
  \label{eq:vanishing_bound}
\end{equation}

When the base of this exponential is less than 1, the gradient shrinks
exponentially in the temporal distance $(t-k)$, making it impossible for the
network to learn that $h_k$ should influence $L_t$ for large $(t-k)$.

\begin{keyinsight}[title=Exponential Gradient Decay]
For a sequence of length $T$, the gradient signal from the last time step that
reaches the first time step has passed through $T$ multiplications by
$\Whh^{\top}$ and $T$ applications of $\tanh'(\cdot)$.  If each factor
contributes a factor of, say, $0.9$, then after 50 steps the gradient magnitude
is $0.9^{50} \approx 0.005$ -- less than $1\%$ of its original strength.  After
100 steps it is $0.9^{100} \approx 2.6 \times 10^{-5}$.
\end{keyinsight}

\subsection{Local Influence Effect}
\label{subsec:local_influence}

The practical consequence of the vanishing gradient is what the instructors
called the \emph{local influence} or \emph{local effect}: each output $\hat{y}_t$
is effectively influenced only by the inputs in its temporal neighbourhood.
Inputs far back in time have negligible influence because the gradient signal
linking them to the loss has vanished.

\begin{figure}[ht]
\centering
\begin{tikzpicture}[
    node distance=0.8cm and 1.6cm,
    cell/.style={draw, rectangle, rounded corners=3pt,
                 minimum width=1.1cm, minimum height=0.8cm,
                 fill=blue!15, thick, font=\small},
    io/.style={draw, circle, minimum size=0.7cm,
               fill=orange!20, thick, font=\small},
    arr/.style={-{Stealth[length=5pt]}, thick},
    fade/.style={-{Stealth[length=5pt]}, thick, draw=gray!40},
    highlight/.style={-{Stealth[length=5pt]}, ultra thick, draw=red!70!black}
  ]
  % Hidden states
  \node[cell] (h0) {$h_0$};
  \node[cell, right=of h0] (h1) {$h_1$};
  \node[cell, right=of h1] (h2) {$h_2$};
  \node[cell, right=of h2] (h3) {$h_3$};
  \node[cell, right=of h3] (h4) {$h_4$};

  % Recurrent arrows: strong near end, fade near start
  \draw[fade]      (h0) -- node[above,font=\tiny]{$\Whh$} (h1);
  \draw[fade]      (h1) -- node[above,font=\tiny]{$\Whh$} (h2);
  \draw[arr]       (h2) -- node[above,font=\tiny]{$\Whh$} (h3);
  \draw[highlight] (h3) -- node[above,font=\tiny]{$\Whh$} (h4);

  % Outputs
  \node[io, above=of h4] (y4) {$\hat{y}_4$};
  \draw[arr] (h4) -- (y4);

  % Gradient back-arrows
  \draw[fade,      bend left=30] (y4) to node[below,font=\tiny,xshift=5pt]{tiny} (h0.south);
  \draw[fade,      bend left=25] (y4) to node[below,font=\tiny]{small} (h1.south);
  \draw[arr,       bend left=20] (y4) to (h2.south);
  \draw[highlight, bend left=10] (y4) to node[below,font=\tiny]{strong} (h3.south);

  % Legend
  \node[right=0.4cm of h4, font=\footnotesize, align=left] (leg) {
    \textcolor{red!70!black}{\rule{0.5cm}{2pt}} strong gradient \\
    \textcolor{black}{\rule{0.5cm}{1pt}} moderate \\
    \textcolor{gray!60}{\rule{0.5cm}{1pt}} vanished
  };

\end{tikzpicture}
\caption{Illustration of the local influence effect.  The gradient signal from
output $\hat{y}_4$ reaches $h_3$ and $h_4$ strongly (red), but effectively
vanishes by the time it reaches $h_0$ and $h_1$ (gray).  Only nearby inputs
influence the current output -- long-range dependencies cannot be learned.}
\label{fig:local_influence}
\end{figure}

\subsection{Two Sources of Vanishing Gradient in RNNs}
\label{subsec:two_sources}

The instructors emphasised that in RNNs there are \emph{two} independent sources
of the vanishing gradient problem, unlike plain deep MLPs which have only one:

\begin{enumerate}[label=\arabic*.]
  \item \textbf{Vertical vanishing gradient:} At any single time step, the
        unrolled RNN is essentially a deep MLP (if the recurrent cell has
        multiple layers).  The conventional backpropagation through these
        vertical layers can suffer from vanishing gradients, exactly as in any
        deep network.
  \item \textbf{Horizontal vanishing gradient (BPTT):} The gradient must also
        travel horizontally across all time steps.  Each step multiplies by
        $\Whh^{\top}$ and $\tanh'(\cdot)$, causing the gradient to vanish
        exponentially in the sequence length.
\end{enumerate}

\begin{warningbox}[title=Compounded Vanishing Gradient in RNNs]
Because RNNs have \emph{both} vertical and horizontal gradient paths, the
vanishing gradient problem is more severe in RNNs than in ordinary deep
networks.  In a deep MLP with $L$ layers, gradients vanish over depth $L$.  In
an RNN with sequence length $T$ and recurrent cell depth $D$, gradients can
vanish over both $D$ (vertical) and $T$ (horizontal), making training on long
sequences extremely difficult.
\end{warningbox}

\subsection{Experimental Illustration}
\label{subsec:experiment}

The instructor shared a student experiment on the IMDB sentiment classification
dataset (movie reviews, binary positive/negative labels) to illustrate the
importance of sequence knowledge.  \Cref{tab:imdb} shows representative results.

\begin{table}[ht]
\centering
\caption{Indicative accuracy on IMDB sentiment classification: feedforward MLP
vs.\ RNN at different sequence lengths.  The MLP processes all words in
parallel (no sequence knowledge); the RNN processes them sequentially.
Note: the instructor cautioned that these numbers are illustrative only.}
\label{tab:imdb}
\renewcommand{\arraystretch}{1.3}
\begin{tabular}{@{}cccc@{}}
\toprule
\textbf{Sequence Length} &
\textbf{MLP Accuracy (\%)} &
\textbf{RNN Accuracy (\%)} &
\textbf{RNN Advantage} \\
\midrule
50  & $\approx 50$ & $> 50$ & Moderate \\
100 & $\approx 50$ & $> 50$ & Moderate \\
\bottomrule
\end{tabular}
\end{table}

The MLP achieves roughly 50\% accuracy (chance level for binary classification)
at any sequence length because it cannot model the sequential context needed to
determine sentiment.  The RNN outperforms the MLP because it accumulates
sequence knowledge step by step.  However, as the sequence length grows, the
RNN's advantage eventually diminishes because long-range gradients vanish and
the network loses track of early context.

%% ══════════════════════════════════════════════
\section{RNN Types and Applications}
\label{sec:rnn_types}
%% ══════════════════════════════════════════════

\subsection{Taxonomy of RNN Configurations}
\label{subsec:rnn_taxonomy}

RNNs are not restricted to the many-to-many configuration shown in the unrolled
diagram.  The input and output sequence lengths are generally independent.  The
lecture discussed the following configurations:

\begin{table}[ht]
\centering
\caption{Common RNN input/output configurations and representative applications.}
\label{tab:rnn_configs}
\renewcommand{\arraystretch}{1.4}
\begin{tabularx}{\textwidth}{@{}lllX@{}}
\toprule
\textbf{Type} & \textbf{Input} & \textbf{Output} & \textbf{Applications} \\
\midrule
One-to-one     & Single vector & Single vector &
  Standard classification (MLP; RNN not needed) \\
Many-to-one    & Sequence      & Single vector &
  Sentiment analysis, document classification \\
One-to-many    & Single vector & Sequence      &
  Image captioning, music generation \\
Many-to-many (sync) & Sequence & Sequence (same length) &
  Named entity recognition, POS tagging \\
Many-to-many (async) & Sequence & Sequence (different length) &
  Machine translation, speech recognition \\
\bottomrule
\end{tabularx}
\end{table}

\begin{keyinsight}[title=Input and Output Lengths are Independent]
A common misconception is that the input and output sequence lengths must be
equal.  They need not be.  In machine translation, an English sentence of 3
words (``I am going'') may translate to a Hindi sentence of 4 words.  The
asymmetric many-to-many case requires an encoder--decoder architecture
(sequence-to-sequence), which is covered in later sessions alongside
Transformers.
\end{keyinsight}

\subsection{Speech Recognition as a Running Example}
\label{subsec:speech_example}

The instructors used speech recognition throughout the session as a concrete
illustration.  In this task:

\begin{itemize}[noitemsep]
  \item \textbf{Input:} A recorded speech signal is processed into a sequence
        of acoustic feature vectors $x_1, x_2, \ldots, x_T$ (e.g., one feature
        vector per 10\,ms frame).  For the phrase ``recurrent neural network'',
        this might produce $T \approx 50$--100 feature vectors.
  \item \textbf{Output:} The three text words $\{y_1, y_2, y_3\}$ = \{recurrent,
        neural, network\}.
  \item \textbf{Length mismatch:} $T_{\text{input}} \approx 50$--100 while
        $T_{\text{output}} = 3$.  A different recording of the same phrase
        might have $T_{\text{input}} = 47$ or $T_{\text{input}} = 52$, but
        $T_{\text{output}}$ remains 3.
\end{itemize}

\begin{example}[title=Word-by-Word Input Processing]
Consider the sentence ``We are in LLM class'' (5 words, vocabulary size $V$).
At each time step one word is fed in as a one-hot vector of length $V$:

\medskip
\begin{center}
\renewcommand{\arraystretch}{1.2}
\begin{tabular}{@{}ccc@{}}
\toprule
\textbf{Time step} & \textbf{Input $x_t$} & \textbf{One-hot index} \\
\midrule
$t=0$ & ``We''    & index of ``We'' in vocabulary \\
$t=1$ & ``are''   & index of ``are'' \\
$t=2$ & ``in''    & index of ``in'' \\
$t=3$ & ``LLM''   & index of ``LLM'' \\
$t=4$ & ``class'' & index of ``class'' \\
\bottomrule
\end{tabular}
\end{center}
\medskip
At each step the hidden state is updated:
$h_t = \tanh(\Wxh\, x_t + \Whh\, h_{t-1})$.
By step $t=4$, $h_4$ encodes contextual information from all five words.
\end{example}

\subsection{Input Representation: One-Hot Encoding}
\label{subsec:one_hot}

Words cannot be fed directly into a neural network; they must be represented
numerically.  The simplest representation is \emph{one-hot encoding}: a vector
of length $V$ (vocabulary size) with a 1 in the position corresponding to the
word and 0 everywhere else.

For a vocabulary $\{$``we'', ``are'', ``LLM'', ``class''$\}$ of size $V=4$:

\begin{align}
  \text{``we''}   &\mapsto [1, 0, 0, 0]^\top \notag \\
  \text{``are''}  &\mapsto [0, 1, 0, 0]^\top \notag \\
  \text{``LLM''}  &\mapsto [0, 0, 1, 0]^\top \notag \\
  \text{``class''}&\mapsto [0, 0, 0, 1]^\top \notag
\end{align}

The vector length $V$ determines the input dimension $d_x$ and hence the number
of columns in $\Wxh$.  The instructor noted that more sophisticated
representations (word embeddings, dense vectors) can be used in place of
one-hot vectors, and would be covered in subsequent sessions.

\subsection{Recurrent CNN (RCNN)}
\label{subsec:rcnn}

A student question at the end of the session asked whether the recurrent
connection concept can be combined with CNNs.  The instructor confirmed this
exists and is called a \emph{Recurrent CNN (RCNN)}, where convolutional layers
are augmented with recurrent connections.  The course does not cover RCNN in
detail, but students were encouraged to explore it independently.

%% ══════════════════════════════════════════════
\section{Solutions to the Vanishing Gradient Problem}
\label{sec:solutions}
%% ══════════════════════════════════════════════

\subsection{Overview of Mitigation Strategies}
\label{subsec:solutions_overview}

Several approaches have been proposed to address or mitigate the vanishing
gradient problem in RNNs:

\begin{enumerate}[label=\arabic*.]
  \item \textbf{Careful weight initialisation:} Initialising $\Whh$ as an
        orthogonal matrix ensures $\|\Whh\| = 1$ initially, delaying (but not
        eliminating) gradient decay.
  \item \textbf{Gradient clipping:} Cap the gradient norm before applying the
        weight update to prevent \emph{exploding} gradients.  Does not address
        \emph{vanishing} gradients directly.
  \item \textbf{ReLU activation:} Unlike $\tanh$, the ReLU function has
        derivative 1 for positive inputs, avoiding the saturation-driven
        component of the vanishing gradient.  However, this introduces the
        dying-ReLU problem and instability in the recurrent path.
  \item \textbf{LSTM (Long Short-Term Memory):} Introduces gating mechanisms
        (forget gate, input gate, output gate) that create direct gradient
        highways through the cell state, largely bypassing the multiplicative
        gradient decay.  Introduced by Hochreiter \& Schmidhuber (1997).
  \item \textbf{GRU (Gated Recurrent Unit):} A simplified variant of LSTM with
        fewer gates (reset gate, update gate), introduced by Cho et al.\ (2014).
        Achieves comparable performance to LSTM with fewer parameters.
\end{enumerate}

\begin{keyinsight}[title=LSTM and GRU as Gradient Highways]
The fundamental innovation of LSTM and GRU is not simply ``more parameters'' but
the creation of \emph{additive} gradient paths.  In a plain RNN the gradient
from step $T$ to step $k$ must pass through $T-k$ \emph{multiplicative}
operations.  In an LSTM, the cell state update includes an \emph{additive}
component (the forget gate times the old cell state plus the input gate times
the new candidate), which allows gradients to flow backward with minimal
attenuation.
\end{keyinsight}

\subsection{Comparison: Plain RNN vs.\ LSTM vs.\ GRU}
\label{subsec:rnn_lstm_gru_compare}

\begin{table}[ht]
\centering
\caption{Comparison of Plain RNN, LSTM, and GRU on key properties.}
\label{tab:rnn_vs_lstm_vs_gru}
\renewcommand{\arraystretch}{1.4}
\begin{tabularx}{\textwidth}{@{}lXXX@{}}
\toprule
\textbf{Property} & \textbf{Plain RNN} & \textbf{LSTM} & \textbf{GRU} \\
\midrule
Hidden state     & $h_t$ only & $h_t$ + cell state $c_t$ & $h_t$ only \\
Gates            & None & Forget, Input, Output (3) & Reset, Update (2) \\
Gradient flow    & Multiplicative chain & Additive cell state path & Additive update path \\
Vanishing gradient & Severe & Largely mitigated & Largely mitigated \\
Parameters       & $O(d_h^2)$ & $O(4 d_h^2)$ & $O(3 d_h^2)$ \\
Training speed   & Fastest & Slowest & Intermediate \\
Long-range deps. & Poor & Good & Good \\
Introduced       & 1986 (Elman) & 1997 (Hochreiter \& Schmidhuber) & 2014 (Cho et al.) \\
\bottomrule
\end{tabularx}
\end{table}

\subsection{Course Roadmap for Sequence Models}
\label{subsec:roadmap}

The session concluded by placing the vanilla RNN in the broader context of the
sequence modelling module:

\begin{center}
\begin{tikzpicture}[
    node distance=0.5cm and 2.0cm,
    box/.style={draw, rectangle, rounded corners=3pt,
                minimum width=2.8cm, minimum height=0.9cm,
                fill=blue!10, thick, font=\small, align=center},
    done/.style={box, fill=green!15},
    future/.style={box, fill=gray!15, dashed},
    arr/.style={-{Stealth[length=6pt]}, thick}
  ]
  \node[done]   (s16) {Session 16\\Sequential data\\motivation};
  \node[done,   right=of s16] (s17) {Session 17\\Plain RNN,\\BPTT, Vanishing~$\nabla$};
  \node[future, right=of s17] (s18) {Sessions 18--20\\LSTM};
  \node[future, right=of s18] (s21) {Sessions 21--22\\GRU};
  \node[future, below=0.8cm of s18, xshift=1.5cm] (s_deep) {Sessions 23+\\Deep \& Bidirectional\\RNN};

  \draw[arr] (s16) -- (s17);
  \draw[arr] (s17) -- (s18);
  \draw[arr] (s18) -- (s21);
  \draw[arr] (s18) -- (s_deep);
  \draw[arr] (s21) -- (s_deep);

\end{tikzpicture}
\end{center}

%% ══════════════════════════════════════════════
\section{Summary and Conclusions}
\label{sec:summary}
%% ══════════════════════════════════════════════

Session~17 provided a deep theoretical foundation for Recurrent Neural Networks.
The following key points were covered:

\begin{enumerate}[label=\arabic*., leftmargin=*]

  \item \textbf{MLP limitations for sequential data.}
        MLPs require fixed-size inputs, cannot handle variable-length sequences
        without wasteful zero-padding, lose order information by processing all
        inputs simultaneously, and learn position-specific rather than
        position-invariant patterns.

  \item \textbf{Parameter sharing as the RNN core innovation.}
        The transition from MLP to RNN is achieved by reusing the same weight
        matrices $\{\Wxh, \Whh, \Why\}$ at every time step.  This single change
        simultaneously addresses variable-length inputs, order preservation,
        reduced parameter count, and statistical generalisation across positions.

  \item \textbf{Forward pass equations.}
        The hidden state update
        $h_t = \tanh(\Wxh\, x_t + \Whh\, h_{t-1} + b_h)$
        is the defining equation of the plain RNN.  Compared to the MLP,
        the only addition is the $\Whh\, h_{t-1}$ term carrying sequence memory.

  \item \textbf{Backpropagation Through Time (BPTT).}
        The total loss $L = \sum_t L_t$ is minimised by computing
        $\nabla_W L = \sum_t \nabla_W L_t$, which requires propagating gradients
        both vertically (through depth) and horizontally (through time steps)
        in the unrolled graph.

  \item \textbf{Vanishing gradient problem.}
        Because the horizontal gradient path involves a product of $(t-k)$
        matrices of the form $\Whh^{\top} \cdot \mathrm{diag}(\tanh'(h_i))$,
        gradients decay exponentially with temporal distance.  This leads to a
        \emph{local influence effect}: outputs are influenced only by nearby
        inputs.

  \item \textbf{Two sources of vanishing gradients.}
        RNNs suffer vanishing gradients both vertically (through deep hidden
        layers) and horizontally (through long sequences via BPTT).  This
        double exposure makes the problem strictly worse in RNNs than in
        standard deep MLPs.

  \item \textbf{LSTM and GRU as solutions.}
        Long Short-Term Memory (LSTM) and Gated Recurrent Unit (GRU)
        architectures were specifically designed to provide additive gradient
        pathways that resist exponential decay, enabling learning of long-range
        dependencies.  These are covered in detail in Sessions~18--22.

\end{enumerate}

\begin{keyinsight}[title=The Central Tension of Plain RNNs]
The plain RNN is theoretically capable of modelling arbitrarily long
dependencies, but is practically unable to learn them because of the vanishing
gradient problem.  The entire purpose of LSTM and GRU is to resolve this
tension by giving the network a mechanism to \emph{choose} which information to
remember and for how long -- something the plain RNN cannot do.
\end{keyinsight}

\vspace{1cm}
\noindent\textbf{Acknowledgements.}  These notes are based on the lecture
delivered by Prof.\ S R M Prasanna and Dr.\ Sunil Saumya at IIIT Dharwad,
January 13, 2026.  The slide deck \texttt{05-DL-RNN} (slides 17--43) and the
session video at \url{https://vimeo.com/1155843264} are the primary sources.
Additional mathematical formulation draws on the standard deep learning
literature, particularly Goodfellow, Bengio \& Courville (2016) and the MIT
6.S191 sequence modelling lecture by Ava Soleimany (2022).

\clearpage

\end{document}
